\section{Multi-Rate Training and Rate Weighting}
\label{sec:multi-rate}

\subsection{Joint Multi-Rate Training}

The model is trained jointly on data from three rate points: r2, r4, and r6.
At each training step, a mini-batch of patches is assembled from the decoded
clouds at all three rates; the rate scalar $r_k = \log_2(\bpp_k)$ is passed
to the FiLM head for each sample $k$.
The total loss is a weighted sum:
\begin{equation}
  \mathcal{L}_{\text{total}} =
    w_2 \mathcal{L}_{\text{CD}}^{(r2)}
    + w_4 \mathcal{L}_{\text{CD}}^{(r4)}
    + w_6 \mathcal{L}_{\text{CD}}^{(r6)},
\end{equation}
where $w_2, w_4, w_6 > 0$ are rate weights.

\subsection{Rate Weight Design}

The choice of weights is non-trivial.
Table~\ref{tab:displacement-distribution} shows that the displacement signal
(mean $\|d_{gt}\|$) decays from r2 to r6 by a factor of $\approx 3$.
This means the Chamfer Distance loss at r6 is intrinsically smaller than at r2,
even for a perfect correction.
Without explicit weighting, the gradient budget would be skewed toward r2
automatically; explicit weights control this balance.

\paragraph{Static weights.}
We use $[w_2, w_4, w_6] = [1.0, 0.3, 0.05]$, assigning most gradient budget
to r2 where the displacement signal is richest and the improvement headroom
is largest.
The rationale: at r6, the geometric correction required is so small that the
gradient contribution is weak regardless of weighting; aggressively upweighting
r6 would not help the r6 predictions while degrading r2 by redirecting gradient
budget away from the most learnable rate.

\paragraph{Why alternative schemes fail.}
Two principled alternatives were evaluated (ablation in
Section~\ref{ssec:ablation-weighting}):
\begin{itemize}
  \item \textit{Kendall uncertainty weighting~\cite{kendall2018multi}}:
    learns weights $w_k \propto 1/\sigma_k^2$ where $\sigma_k$ is a per-task
    uncertainty parameter.
    In practice, the highest-loss task (r2) is interpreted as highest-uncertainty
    and \emph{downweighted}: weights collapse to the inverted ordering
    $[w_2, w_4, w_6] \approx [0.1, 0.3, 0.6]$ within 5 epochs.
    This is the opposite of the correct inductive bias.

  \item \textit{Dynamic detached weighting}:
    weights proportional to current loss magnitude, normalised and clamped.
    Converges quickly to a nearly static ratio $\approx [1.27, 0.91, 0.82]$
    because all rates improve at proportional rates (the shared backbone
    prevents decoupled learning dynamics between rates).
    Marginally worse than static weights.
\end{itemize}

The conclusion is that the static weights $[1.0, 0.3, 0.05]$ encode the
correct inductive prior --- most budget to r2 --- and no adaptive scheme
tested can recover a better solution from this regime.

% ============================================================
